\section*{Приложение}
\addcontentsline{toc}{section}{Приложение}
\label{sec:Apendix} \index{Apendix}
\hypertarget{sec:Appendix_hypertarget}{}

%Здесь необходимо написать приложение, которое вы должны придумать самостоятельно.

В приложении представлены листинги программного кода, на которые имеются ссылки в основном тексте работы.

%===============lst:vm_replica_flags===============

\begin{lstlisting}[style=diff]
 /* include/linux/mm.h */

 #ifdef CONFIG_ARCH_USES_HIGH_VMA_FLAGS
 #define VM_HIGH_ARCH_BIT_0	32
 ...
 #define VM_HIGH_ARCH_BIT_5	37
<+>+#define VM_HIGH_ARCH_BIT_6	39</+>
<+>+#define VM_HIGH_ARCH_BIT_7	40</+>
 #define VM_HIGH_ARCH_0	BIT(VM_HIGH_ARCH_BIT_0)
 ...
 #define VM_HIGH_ARCH_5	BIT(VM_HIGH_ARCH_BIT_5)
<+>+#define VM_HIGH_ARCH_6	BIT(VM_HIGH_ARCH_BIT_6)</+>
<+>+#define VM_HIGH_ARCH_7	BIT(VM_HIGH_ARCH_BIT_7)</+>
 #endif /* CONFIG_ARCH_USES_HIGH_VMA_FLAGS */
 ...
<+>+#ifdef CONFIG_USER_REPLICATION</+>
<+>+/* Page tables for this vma should be replicated during page faults */</+>
<+>+#define VM_REPLICA_INIT	VM_HIGH_ARCH_6</+>
<+>+/*</+>
<+>+ * Phys memory of this vma might/should be replicated.</+>
<+>+ * If this flag is set, VM_REPLICA_INIT also must be set.</+>
<+>+ */</+>
<+>+#define VM_REPLICA_COMMIT	VM_HIGH_ARCH_7</+>
<+>+#endif /* CONFIG_USER_REPLICATION */</+>
\end{lstlisting}
\begin{lstlisting}[
    style=normal,
    caption={Флаги репликации для виртуальных областей памяти},
    label={lst:vm_replica_flags}
]
/* include/linux/numa_user_replication.h */

#ifdef CONFIG_USER_REPLICATION
...
static inline bool numa_is_vma_replicant(struct vm_area_struct *vma) {
	return vma->vm_flags & VM_REPLICA_INIT;
}

static inline bool vma_has_replicas(struct vm_area_struct *vma) {
	return vma->vm_flags & VM_REPLICA_COMMIT;
}
...
#else
...
static inline bool numa_is_vma_replicant(struct vm_area_struct *vma) {
	return 0;
}

static inline bool vma_has_replicas(struct vm_area_struct *vma) {
	return 0;
}
...
#endif /* CONFIG_USER_REPLICATION */
\end{lstlisting}
\vspace{0.5cm}

%===============lst:mm_struct===============

\begin{lstlisting}[style=diff]
 /* include/linux/mm_types.h */

<+>+enum table_replication_policy_t {</+>
<+>+	/* no table replication */</+>
<+>+	TABLE_REPLICATION_NONE = 0,</+>
<+>+	/* replicate minimal amount to support RO-data replication */</+>
<+>+	TABLE_REPLICATION_MINIMAL,</+>
<+>+	/* replicate all tables */</+>
<+>+	TABLE_REPLICATION_ALL</+>
<+>+};</+>
<+>+</+>
<+>+enum data_replication_policy_t {</+>
<+>+	/* no data replicas at all */</+>
<+>+	DATA_REPLICATION_NONE = 0,</+>
<+>+	/* replicas are created on demand via numa balancer */</+>
<+>+	DATA_REPLICATION_ON_DEMAND,</+>
<+>+	/* all ro-data always replicated */</+>
<+>+	DATA_REPLICATION_ALL</+>
<+>+};</+>
<+>+</+>
<+>+struct numa_replication_policy_t {</+>
<+>+	enum table_replication_policy_t table_policy:3;</+>
<+>+	enum data_replication_policy_t data_policy:3;</+>
<+>+};</+>
<+>+</+>
 ...
 struct mm_struct {
     ...
<+>+#ifdef CONFIG_USER_REPLICATION</+>
<+>+    struct numa_replication_policy_t replication_policy;</+>
<+>+#endif</+>
 };
\end{lstlisting}
\begin{lstlisting}[
    style=normal,
    caption={Политики NUMA-репликации},
    label={lst:mm_struct}
]
/* include/linux/numa_user_replication.h */

#ifdef CONFIG_USER_REPLICATION
...
static inline enum table_replication_policy_t
get_table_replication_policy(struct mm_struct *mm) {
	return mm->replication_policy.table_policy;
}

static inline void
set_table_replication_policy(struct mm_struct *mm,
	enum table_replication_policy_t table_pol) {
	mm->replication_policy.table_policy = table_pol;
}

static inline enum data_replication_policy_t
get_data_replication_policy(struct mm_struct *mm) {
	return mm->replication_policy.data_policy;
}

static inline void
set_data_replication_policy(struct mm_struct *mm,
	enum data_replication_policy_t data_pol) {
	mm->replication_policy.data_policy = data_pol;
}
...
#endif /* CONFIG_USER_REPLICATION */
\end{lstlisting}
\vspace{0.5cm}

%===============lst:vmflags_modify_replica===============

\begin{lstlisting}[style=normal]
/* include/linux/numa_user_replication.h */

#ifdef CONFIG_USER_REPLICATION
...
/* true if vma with @vm_flags is suitable for data replication */
static inline bool __vmflags_replica_candidate(vm_flags_t vm_flags) {
	return (vm_flags & VM_ACCESS_FLAGS) &&
        !(vm_flags & (VM_WRITE | VM_SHARED));
}

/* true if vma with @vm_flags is ready for VM_REPLICA_COMMIT flag, i.e.
 * it is suitable for data replication and VM_REPLICA_INIT is already
 * set */
static inline bool vmflags_might_be_replicated(vm_flags_t vm_flags) {
	return (vm_flags & VM_REPLICA_INIT) && (vm_flags & VM_ACCESS_FLAGS) &&
        !(vm_flags & (VM_WRITE | VM_SHARED));
}

/* set replication flags of @vm_flags according to @mm's policies */
static inline void vmflags_modify_replica(
    struct mm_struct *mm, unsigned long *vm_flags) {
	*vm_flags &= ~(VM_REPLICA_INIT | VM_REPLICA_COMMIT);

	if ((get_table_replication_policy(mm) == TABLE_REPLICATION_ALL) ||
		(get_table_replication_policy(mm) == TABLE_REPLICATION_MINIMAL &&
			__vmflags_replica_candidate(*vm_flags))) {
		*vm_flags |= VM_REPLICA_INIT;
		if (get_data_replication_policy(mm) != DATA_REPLICATION_NONE &&
				vmflags_might_be_replicated(*vm_flags))
			*vm_flags |= VM_REPLICA_COMMIT;
	}
}
...
#endif /* CONFIG_USER_REPLICATION */
\end{lstlisting}
\begin{lstlisting}[
    style=diff,
    caption={Установка флагов репликации новым/модифицированным областям},
    label={lst:vmflags_modify_replica}
]
 /* mm/mmap.c */

 unsigned long do_mmap(struct file *file, unsigned long addr,
             unsigned long len, unsigned long prot,
 	         unsigned long flags, vm_flags_t vm_flags,
 			 unsigned long pgoff, unsigned long *populate,
 			 struct list_head *uf)
 {
     struct mm_struct *mm = current->mm;
     ...
<+>+#ifdef CONFIG_USER_REPLICATION</+>
<+>+    vmflags_modify_replica(mm, &vm_flags);</+>
<+>+#endif</+>
     addr = mmap_region(file, addr, len, vm_flags, pgoff, uf);
     ...
     return addr;
 }

 /* mm/mprotect.c */

 static int do_mprotect_pkey(unsigned long start, size_t len,
         unsigned long prot, int pkey)
 {
     unsigned long end;
     struct vm_area_struct *vma;
     struct vma_iterator vmi;
     ...
     for_each_vma_range(vmi, vma, end) {
         unsigned long mask_off_old_flags, newflags;
         int new_vma_pkey;
         ...
         newflags = calc_vm_prot_bits(prot, new_vma_pkey);
         newflags |= (vma->vm_flags & ~mask_off_old_flags);
<+>+#ifdef CONFIG_USER_REPLICATION</+>
<+>+        if (vma->vm_mm)</+>
<+>+            vmflags_modify_replica(vma->vm_mm, &newflags);</+>
<+>+#endif</+>
         ...
     }
     ...
 }
\end{lstlisting}
\vspace{0.5cm}

%===============lst:numa_mm_handle_table_replication===============

\begin{lstlisting}[
    style=normal,
    caption={Установка флага репликации таблиц трансляций при смене политики},
    label={lst:numa_mm_handle_table_replication}
]
/* mm/numa_user_replication.c */

/* called when table replication policy goes from All to Minimal */
static void dereplicate_all2min_pgtables(struct mm_struct *mm) {
	struct vm_area_struct *vma;
    ...
	do {
		unsigned long addr = vma->vm_start;
		struct vm_area_struct *next;
        ...
        /* clear VM_REPLICA_INIT and dereplicate
         * pgtables for not replica candidate vmas */
		if (numa_is_vma_replicant(vma) && !vma_might_be_replicated(vma)) {
			vm_flags_clear(vma, VM_REPLICA_INIT);
			...
			dereplicate_pgtables_pgd_range(&tlb, addr, vma->vm_end, ...);
		}
		vma = next;
	} while (vma);
    ...
}

static void numa_mm_table_replication_minimal(struct mm_struct *mm) {
	switch (get_table_replication_policy(mm)) {
    /* executed when table replication policy goes from None to Minimal */
	case TABLE_REPLICATION_NONE: {
		struct vm_area_struct *vma;
		MA_STATE(mas, &mm->mm_mt, 0, 0);

		mas_for_each(&mas, vma, ULONG_MAX) {
            /* set VM_REPLICA_INIT and replicate
             * pgtables for replica candidate vmas */
			if (vma_replica_candidate(vma)) {
				vm_flags_set(vma, VM_REPLICA_INIT);
				replicate_pgtables_vma(vma);
			}
		}
		break;
	}
	case TABLE_REPLICATION_MINIMAL:
		return;
	case TABLE_REPLICATION_ALL:
		dereplicate_all2min_pgtables(mm);
		break;
	}

	set_table_replication_policy(mm, TABLE_REPLICATION_MINIMAL);
}
\end{lstlisting}
\vspace{0.5cm}

%===============lst:numa_mm_handle_data_replication===============

\begin{lstlisting}[
    style=normal,
    caption={Установка флага репликации данных при смене политики},
    label={lst:numa_mm_handle_data_replication}
]
/* mm/numa_user_replication.c */

static bool numa_mm_data_replication_none(struct mm_struct *mm) {
	switch (get_data_replication_policy(mm)) {
	case DATA_REPLICATION_NONE:
		return true;
    /* executed when data replication policy
     * goes from Lazy/All to None */
	case DATA_REPLICATION_ON_DEMAND:
	case DATA_REPLICATION_ALL: {
		struct vm_area_struct *vma;
		MA_STATE(mas, &mm->mm_mt, 0, 0);

		mas_for_each(&mas, vma, ULONG_MAX) {
            /* deduplicate data and clear VM_REPLICA_COMMIT
             * for all has_replicas vmas */
			if (vma_has_replicas(vma)) {
				vma_start_write(vma);
				if (phys_deduplicate(vma,
						vma->vm_start, vma->vm_end - vma->vm_start))
					return false;
				__vm_flags_mod(vma, VM_NONE, VM_REPLICA_COMMIT);
			}
		}
		break;
	}
	}

	set_data_replication_policy(mm, DATA_REPLICATION_NONE);
	return true;
}

static bool numa_mm_data_replication_all(struct mm_struct *mm) {
	switch (get_data_replication_policy(mm)) {
    /* executed when data replication policy
     * goes from Lazy/None to All */
	case DATA_REPLICATION_NONE:
	case DATA_REPLICATION_ON_DEMAND: {
		struct vm_area_struct *vma;
		MA_STATE(mas, &mm->mm_mt, 0, 0);

		mas_for_each(&mas, vma, ULONG_MAX) {
            /* set VM_REPLICA_COMMIT and duplicate data
             * for all replica candidate vmas */
			if (vma_might_be_replicated(vma))
				vm_flags_set(vma, VM_REPLICA_COMMIT);
			if (vma_has_replicas(vma))
				phys_duplicate(vma,
                    vma->vm_start, vma->vm_end - vma->vm_start);
		}
		break;
	}
	case DATA_REPLICATION_ALL:
		return true;
	}

	set_data_replication_policy(mm, DATA_REPLICATION_ALL);
	return true;
}
\end{lstlisting}
\vspace{0.5cm}

%===============lst:PG_replicated===============

\begin{lstlisting}[
	style=diff,
	caption={Флаг для реплицированных физических страниц},
	label={lst:PG_replicated}
]
 /* include/linux/page-flags.h */

 enum pageflags {
	 PG_locked,
	 ...
<+>+#ifdef CONFIG_USER_REPLICATION</+>
<+>+	 PG_replicated,</+>
<+>+#endif</+>
 	 __NR_PAGEFLAGS,
	 ...
 };
 ...
<+>+#ifdef CONFIG_USER_REPLICATION</+>
<+>+PAGEFLAG(Replicated, replicated, PF_ANY)</+>
<+>+#else</+>
<+>+PAGEFLAG_FALSE(Replicated, replicated)</+>
<+>+#endif</+>
 ...
<+>+#ifdef CONFIG_USER_REPLICATION</+>
<+>+#define __PG_REPLICATED		(1UL << PG_replicated)</+>
<+>+#else</+>
<+>+#define __PG_REPLICATED		0</+>
<+>+#endif</+>
<+>+</+>
 /*
  * Flags checked when a page is freed.  Pages being freed should not have
  * these flags set.  If they are, there is a problem.
  */
 #define PAGE_FLAGS_CHECK_AT_FREE					\
 	 1UL << PG_private    | 1UL << PG_private_2 |	\
 	 1UL << PG_writeback  | 1UL << PG_reserved  |	\
 	 1UL << PG_slab       | 1UL << PG_active    |	\
<->-	 1UL << PG_unevictable| __PG_MLOCKED | LRU_GEN_MASK)</->
<+>+	 1UL << PG_unevictable| __PG_REPLICATED | __PG_MLOCKED | LRU_GEN_MASK)</+>
\end{lstlisting}
\vspace{0.5cm}

%===============lst:master_table===============

\begin{lstlisting}[
	style=diff,
	caption={Использование блокировки \textit{мастер}-таблицы},
	label={lst:master_table}
]
 /* include/linux/mm_types.h */

 struct ptdesc {
	 ...
	 union {
		 struct mm_struct *pt_mm; /* Used for x86 pgds. */
		 atomic_t pt_frag_refcount; /* Powerpc only. */
<+>+#ifdef CONFIG_USER_REPLICATION</+>
<+>+		 struct ptdesc *master_table;</+>
<+>+#endif</+>
	 };

	 union {
 #if ALLOC_SPLIT_PTLOCKS
		 spinlock_t *ptl;
 #else
		 spinlock_t ptl;
 #endif
	 };
	 ...
 };

 /* include/linux/mm.h */

 #if USE_SPLIT_PTE_PTLOCKS
 #if ALLOC_SPLIT_PTLOCKS
 ...
 bool ptlock_alloc(struct ptdesc *ptdesc);
 void ptlock_free(struct ptdesc *ptdesc);

 static inline spinlock_t *ptlock_ptr(struct ptdesc *ptdesc) {
<+>+#ifdef CONFIG_USER_REPLICATION</+>
<+>+	 return ptdesc->master_table->ptl;</+>
<+>+#else</+>
	 return ptdesc->ptl;
<+>+#endif</+>
 }
 #else /* ALLOC_SPLIT_PTLOCKS */
 ...
 static inline bool ptlock_alloc(struct ptdesc *ptdesc) {
<+>+#ifdef CONFIG_USER_REPLICATION</+>
<+>+	 ptdesc->master_table = ptdesc;</+>
<+>+#endif</+>
	 return true;
 }

 static inline void ptlock_free(struct ptdesc *ptdesc) {
<+>+#ifdef CONFIG_USER_REPLICATION</+>
<+>+	 ptdesc->master_table = NULL;</+>
<+>+#endif</+>
 }

 static inline spinlock_t *ptlock_ptr(struct ptdesc *ptdesc) {
<+>+#ifdef CONFIG_USER_REPLICATION</+>
<+>+	 return &ptdesc->master_table->ptl;</+>
<+>+#else</+>
	 return &ptdesc->ptl;
<+>+#endif</+>
 }
 #endif /* ALLOC_SPLIT_PTLOCKS */
 ...
 #endif /* USE_SPLIT_PTE_PTLOCKS */
\end{lstlisting}
\vspace{0.5cm}

%===============lst:replica_list===============

\begin{lstlisting}[style=diff]
 /* include/linux/llist.h */

 struct llist_head {
	 struct llist_node *first;
 };
 
 struct llist_node {
	 struct llist_node *next;
 };

 /* include/linux/mm_types.h */

 struct ptdesc {
	 ...
<->-	 unsigned long __page_mapping; /* Unused for page tables. */</->
<+>+	 union {</+>
<+>+		 unsigned long __page_mapping; /* Unused for page tables. */</+>
<+>+		 struct llist_head replica_list_head;</+>
<+>+		 struct llist_node replica_list_node;</+>
<+>+	 };</+>
	 ...
 };
\end{lstlisting}
\begin{lstlisting}[
	style=normal,
	caption={Циклический односвязный список таблиц-реплик},
	label={lst:replica_list}
]
/* include/linux/numa_user_replication.h */

#ifdef CONFIG_USER_REPLICATION
...
/**
 * @pos:	struct ptdesc* of current replica
 * @table:	current table entry to write (virtual address)
 * @head_table:	table entry to start with, will be a part of this loop
 * @nid:	node id of current pgtable
 * @offset:	offset of current table entry in table page in bytes
 * @start:	boolean value for tmp storage
 */
#define for_each_pgtable(pos, table, head_table, nid, offset, start)	\
	for (pos = llist_entry(									\
			&virt_to_ptdesc(head_table)->replica_list_node,	\
		typeof(*pos), replica_list_node),					\
	    start = true, nid = page_to_nid(ptdesc_page(pos)),	\
	    offset = offset_in_table(head_table),				\
		table = (typeof(table))get_table_ptr(pos, offset);	\
	    pos != virt_to_ptdesc(head_table) || start;			\
	    pos = llist_entry((pos)->replica_list_node.next,	\
			typeof(*pos), replica_list_node),				\
	    table = (typeof(table))get_table_ptr(pos, offset),	\
	    nid = page_to_nid(ptdesc_page(pos)), start = false)

/**
 * @pos:	struct ptdesc* of current replica
 * @table:	current table entry to write (virtual address)
 * @head_table:	table entry to start with, will not be a part of this loop
 * @offset:	offset of current table entry in table page in bytes
 */
#define for_each_pgtable_replica(pos, table, head_table, offset)	\
	for (pos = llist_entry(										\
			virt_to_ptdesc(head_table)->replica_list_node.next,	\
		typeof(*pos), replica_list_node),					\
	    offset = offset_in_table(head_table),				\
		table = (typeof(table))get_table_ptr(pos, offset);	\
	    pos != virt_to_ptdesc(head_table);					\
	    pos = llist_entry((pos)->replica_list_node.next,	\
			typeof(*pos), replica_list_node),				\
	    table = (typeof(table))get_table_ptr(pos, offset))
...
#endif /* CONFIG_USER_REPLICATION */
\end{lstlisting}
\vspace{0.5cm}

%===============lst:pgd_numa===============

\begin{lstlisting}[style=diff]
 /* include/linux/mm_types.h */

 struct mm_struct {
	 struct {
		 ...
		 pgd_t *pgd;
		 ...
	 };
<+>+#ifdef CONFIG_USER_REPLICATION</+>
<+>+	 pgd_t **pgd_numa;</+>
<+>+#endif</+>
	 ...
 };
\end{lstlisting}
\begin{lstlisting}[
	style=normal,
	caption={Указатели на реплицированные PGD-таблицы доступны через \texttt{mm\_struct}},
	label={lst:pgd_numa}
]
/* include/linux/numa_user_replication.h */

#ifdef CONFIG_USER_REPLICATION
...
#define this_node_pgd(mm) ((mm)->pgd_numa[numa_node_id()])
...
#else
...
#define this_node_pgd(mm) ((mm)->pgd)
...
#endif /* CONFIG_USER_REPLICATION */
\end{lstlisting}
\vspace{0.5cm}

%===============lst:handle_mm_fault===============

\begin{lstlisting}[
	style=diff,
	caption={Спуск в иерархии таблиц трансляций в процессе обработки ошибки страниц},
	label={lst:handle_mm_fault}
]
 /* mm/memory.c */

 static vm_fault_t __handle_mm_fault(struct vm_area_struct *vma,
	 unsigned long address, unsigned int flags)
 {
	 struct vm_fault vmf = {...};
	 struct mm_struct *mm = vma->vm_mm;
	 pgd_t *pgd;
	 p4d_t *p4d;
	 ...
<->-	 pgd = pgd_offset(mm, address);</->
<->-	 p4d = p4d_alloc(mm, pgd, address);</->
<+>+	 pgd = fault_pgd_offset(&vmf, address);</+>
<+>+	 p4d = fault_p4d_alloc(&vmf, mm, pgd, address);</+>
	 ...
<->-	 vmf.pud = pud_alloc(mm, p4d, address);</->
<+>+	 vmf.pud = fault_pud_alloc(&vmf, mm, p4d, address);</+>
	 ...
<->-	 vmf.pmd = pmd_alloc(mm, vmf.pud, address);</->
<+>+	 vmf.pmd = fault_pmd_alloc(&vmf, mm, vmf.pud, address);</+>
	 ...
<+>+	 if (fault_pte_alloc(&vmf))</+>
<+>+		 return VM_FAULT_OOM;</+>
	 return handle_pte_fault(&vmf);
 }
\end{lstlisting}
\vspace{0.5cm}

%===============lst:fault_pxx_alloc===============

\begin{lstlisting}[
	style=normal,
	caption={Модифицированные функции для спуска, аллокации и репликации таблиц трансляций},
	label={lst:fault_pxx_alloc}
]
/* include/linux/numa_user_replication.h */

#ifdef CONFIG_USER_REPLICATION
...
pgd_t *fault_pgd_offset(struct vm_fault *vmf, unsigned long address) {
	vmf->pgd = pgd_offset_pgd(this_node_pgd(vmf->vma->vm_mm), address);
	return vmf->pgd;
}

p4d_t *fault_p4d_alloc(struct vm_fault *vmf, struct mm_struct *mm,
	pgd_t *pgd, unsigned long address) {
	if (replication_path_pgd(vmf)) {
		if (replication_handle_p4d_fault(vmf))
			return NULL;
	} else
		vmf->p4d = p4d_alloc(mm, pgd, address);
	return vmf->p4d;
}

pud_t *fault_pud_alloc(struct vm_fault *vmf, struct mm_struct *mm,
	p4d_t *p4d, unsigned long address) {
	if (replication_path_p4d(vmf)) {
		if (replication_handle_pud_fault(vmf))
			return NULL;
	} else
		vmf->pud = pud_alloc(mm, p4d, address);
	return vmf->pud;
}

pmd_t *fault_pmd_alloc(struct vm_fault *vmf, struct mm_struct *mm,
	pud_t *pud, unsigned long address) {
	if (replication_path_pud(vmf)) {
		if (replication_handle_pmd_fault(vmf))
			return NULL;
	} else
		vmf->pmd = pmd_alloc(mm, pud, address);
	return vmf->pmd;
}

int fault_pte_alloc(struct vm_fault *vmf) {
	if (replication_path_pmd(vmf))
		return replication_handle_pte_fault(vmf);
	return 0;
}
...
#else
...
static inline pgd_t *fault_pgd_offset(struct vm_fault *vmf,
	unsigned long address) {
	return pgd_offset(vmf->vma->vm_mm, address);
}

static inline p4d_t *fault_p4d_alloc(struct vm_fault *vmf,
	struct mm_struct *mm, pgd_t *pgd, unsigned long address) {
	return p4d_alloc(mm, pgd, address);
}

static inline pud_t *fault_pud_alloc(struct vm_fault *vmf,
	struct mm_struct *mm, p4d_t *p4d, unsigned long address) {
	return pud_alloc(mm, p4d, address);
}

static inline pmd_t *fault_pmd_alloc(struct vm_fault *vmf,
	struct mm_struct *mm, pud_t *pud, unsigned long address) {
	return pmd_alloc(mm, pud, address);
}

static inline int fault_pte_alloc(struct vm_fault *vmf) {
	return 0;
}
...
#endif /* CONFIG_USER_REPLICATION */
\end{lstlisting}
\vspace{0.5cm}

%===============lst:replication_handle_pte_fault===============

\begin{lstlisting}[
	style=normal,
	caption={Репликация PTE-таблиц},
	label={lst:replication_handle_pte_fault}
]
/* mm/numa_user_replication.c */

static vm_fault_t replication_handle_pte_fault(struct vm_fault *vmf)
{
	int ret;
	struct mm_struct *mm = vmf->vma->vm_mm;
	struct ptdesc *pte_tables[MAX_NUMNODES];
	spinlock_t *ptl;

	/* true if PTE-level is absent or replicated */
	if (replicated_pte_level(vmf)) {
		if (!pmd_none(*vmf->pmd))
			/* PTE-level is not absent, therefore it is replicated.
			 * Nothing need to do, returning. */
			return 0;

		/* PTE-level is absent, we need to create and replicate it */

		/* Allocate and prepare replicated PTE-tables.
		 * Preparation involves:
		 * 1. setting PG_replicated flag;
		 * 2. setting master-table for tables-replicas;
		 * 3. connecting tables into cyclic linked list.
		 */
		ret = prepare_replicated_pte_tables(NUMA_NO_NODE, pte_tables, mm);
		if (ret)
			return VM_FAULT_OOM;
		/* acquire PMD-lock before modifying parent PMD-tables */
		ptl = pmd_lock(mm, vmf->pmd);
		if (unlikely(pmd_present(*vmf->pmd))) {
			/* Someone else has already replicated this level.
			 * Nothing else need to do, releasing and returning. */
			spin_unlock(ptl);
			release_replicated_pte_tables(NUMA_NO_NODE, pte_tables, mm);
		} else {
			/* Populate PMD-entries to point to replicated PTE-tables. */
			sync_replicated_pte_tables(NUMA_NO_NODE, pte_tables,
				vmf->pmd, mm);
			spin_unlock(ptl);
		}
	} else {
		/* the only PTE-table exists, we need to replicate it */

		struct page *table_page = pmd_pgtable(*(vmf->pmd));
		int pte_node = page_to_nid(table_page);

		pte_tables[pte_node] = page_ptdesc(table_page);
		/* Allocate remaining PTE-tables. */
		ret = prepare_replicated_pte_tables(pte_node, pte_tables, mm);
		if (ret)
			return VM_FAULT_OOM;
		/* acquire PMD-lock before modifying parent PMD-tables */
		ptl = pmd_lock(mm, vmf->pmd);
		if (unlikely(PageReplicated(table_page))) {
			/* Someone else has already replicated this level.
			 * Nothing else need to do, releasing and returning. */
			spin_unlock(ptl);
			release_replicated_pte_tables(pte_node, pte_tables, mm);
		} else {
			/* Acquire PTE-lock of existing PTE-table.
			 * Prepare replicated PTE-tables, it involves:
			 * 1. setting PG_replicated flag;
			 * 2. setting master-table to point on existing PTE-table;
			 * 3. connecting tables into cyclic linked list;
			 * 4. copy the contents of existing PTE-table to others.
			 * Populate PMD-entries to point to replicated PTE-tables.
			 */
			sync_replicated_pte_tables(pte_node, pte_tables,
				vmf->pmd, mm);
			spin_unlock(ptl);
		}
	}

	return 0;
}
\end{lstlisting}
\vspace{0.5cm}

%===============lst:data_replication_all===============

\begin{lstlisting}[
	style=diff,
	caption={Репликация данных при установленной политике \texttt{All}},
	label={lst:data_replication_all}
]
 /* mm/memory.c */

 static vm_fault_t do_pte_missing(struct vm_fault *vmf)
 {
<+>+	 vm_fault_t ret;</+>
	 ...
	 if (vma_is_anonymous(vmf->vma))
<->-		 return do_anonymous_page(vmf);</->
<+>+		 ret = do_anonymous_page(vmf);</+>
	 else
<->-		 return do_fault(vmf);</->
<+>+		 ret = do_fault(vmf);</+>
<+>+</+>
<+>+#ifdef CONFIG_USER_REPLICATION</+>
<+>+	 if (get_data_replication_policy(vmf->vma->vm_mm) ==</+>
<+>+		 DATA_REPLICATION_ALL) {</+>
<+>+		 if (vma_has_replicas(vmf->vma))</+>
<+>+			 numa_replicate_page(vmf);</+>
<+>+	 }</+>
<+>+#endif</+>
<+>+	 return ret;</+>
 }

 /* mm/mprotect.c */

 int mprotect_fixup(struct vma_iterator *vmi, struct mmu_gather *tlb,
	 struct vm_area_struct *vma, struct vm_area_struct **pprev,
	 unsigned long start, unsigned long end, unsigned long newflags)
 {
	 struct mm_struct *mm = vma->vm_mm;
	 unsigned int mm_cp_flags = 0;
	 ...
	 vm_flags_reset(vma, newflags);
	 if (vma_wants_manual_pte_write_upgrade(vma))
		 mm_cp_flags |= MM_CP_TRY_CHANGE_WRITABLE;
	 vma_set_page_prot(vma);

	 change_protection(tlb, vma, start, end, mm_cp_flags);

<+>+#ifdef CONFIG_USER_REPLICATION</+>
<+>+	 if (get_data_replication_policy(mm) == DATA_REPLICATION_ALL) {</+>
<+>+		 if (vma_has_replicas(vma)) {</+>
<+>+			 replicate_pgtables_vma(vma);</+>
<+>+			 phys_duplicate(vma, start, end - start);</+>
<+>+		 }</+>
<+>+	 }</+>
<+>+#endif</+>
	 ...
 }
\end{lstlisting}
\vspace{0.5cm}

%===============lst:NUMA_balancer===============

\begin{lstlisting}[
	style=diff,
	caption={Модифицированная фаза сканирования NUMA-балансировщика},
	label={lst:NUMA_balancer}
]
 /* kernel/sched/fair.c */

 static void task_numa_work(struct callback_head *work)
 {
	 struct vm_area_struct *vma;
	 struct vma_iterator vmi;
	 ...
	 for (; vma; vma = vma_next(&vmi)) {
<->-		 if (!vma_migratable(vma) || !vma_policy_mof(vma) ||</->
<+>+		 if (!vma_migratable(vma) || (!vma_policy_mof(vma) &&</+>
<+>+			 !vma_has_replicas(vma)) ||</+>
			 is_vm_hugetlb_page(vma) || (vma->vm_flags & VM_MIXEDMAP))
			 continue;

		 /* Avoid hinting faults in read-only file-backed mappings
		  * as migrating the pages will be of marginal benefit. */
<+>+</+>
<+>+		 /* Do not avoid hinting faults in read-only file-backed mappings</+>
<+>+		  * as replicating the pages should be of significant benefit. */</+>
<->-		 if (!vma->vm_mm || (vma->vm_file &&</->
<->-			 (vma->vm_flags & (VM_READ|VM_WRITE)) == (VM_READ)))</->
<+>+		 if ((!vma->vm_mm || (vma->vm_file &&</+>
<+>+			 (vma->vm_flags & (VM_READ|VM_WRITE)) == (VM_READ))) &&</+>
<+>+			 !vma_has_replicas(vma))</+>
			 continue;
		 ...
	 }
	 ...
 }

 /* mm/mprotect.c */

 static long change_pte_range(struct mmu_gather *tlb,
	 struct vm_area_struct *vma, pmd_t *pmd, unsigned long addr,
	 unsigned long end, pgprot_t newprot, unsigned long cp_flags)
 {
	 pte_t *pte, oldpte;
	 bool prot_numa = cp_flags & MM_CP_PROT_NUMA;
	 long pages = 0;
	 ...
	 do {
		 oldpte = ptep_get(pte);
		 if (pte_present(oldpte)) {
			 ...
			 if (prot_numa) {
				 struct page *page;
				 ...
<+>+				 /* skip already replicated memory */</+>
<+>+				 if (PageReplicated(page))</+>
<+>+				 	 return pages;</+>
				 ...
			 }
			 ...
		 }
	 } while (pte++, addr += PAGE_SIZE, addr != end);
	 ...
	 return pages;
 }
\end{lstlisting}
\vspace{0.5cm}

%===============lst:do_numa_page===============

\begin{lstlisting}[
	style=diff,
	caption={Модифицированный обработчик \textit{NUMA hinting faults}},
	label={lst:do_numa_page}
]
 static vm_fault_t do_numa_page(struct vm_fault *vmf)
 {
	 struct vm_area_struct *vma = vmf->vma;
	 struct page *page = NULL;
	 int page_nid = NUMA_NO_NODE, target_nid;
	 ...
	 target_nid = numa_migrate_prep(
		 page, vma, vmf->address, page_nid, &flags);
<->-	 if (target_nid == NUMA_NO_NODE) {</->
<+>+	 if (target_nid == NUMA_NO_NODE && !vma_has_replicas(vma)) {</+>
		 put_page(page);
		 goto out_map;
	 }
	 ...
<+>+#ifdef CONFIG_USER_REPLICATION</+>
<+>+	 if (vma_has_replicas(vma)) {</+>
<+>+		 /* drop the reference count that was</+>
<+>+		  * elevated in numa_migrate_prep() */</+>
<+>+		 put_page(folio);</+>
<+>+</+>
<+>+		 if (numa_replicate_page(vmf)) {</+>
<+>+			 /* replication success */</+>
<+>+			 ...</+>
<+>+			 return 0;</+>
<+>+		 }</+>
<+>+		 /* replication failed */</+>
<+>+</+>
<+>+		 /* forbid migration if the condition from</+>
<+>+		  * kernel/sched/fair.c:task_numa_work()</+>
<+>+		  * is not satisfied */</+>
<+>+		 if (vma->vm_file &&</+>
<+>+			 ((vma->vm_flags & (VM_READ|VM_WRITE)) == VM_READ)) {</+>
<+>+			 ...</+>
<+>+			 goto out_map;</+>
<+>+		 }</+>
<+>+</+>
<+>+		 /* forbid migration if the condition from</+>
<+>+		  * beginning of the function is not satisfied */</+>
<+>+		 if (target_nid == NUMA_NO_NODE) {</+>
<+>+			 ...</+>
<+>+			 goto out_map;</+>
<+>+		 }</+>
<+>+	 }</+>
<+>+	 /* fallback to migration */</+>
<+>+#endif</+>
	 ...
 }
\end{lstlisting}
\vspace{0.5cm}

%===============lst:dereplicate_pgtables_pte_range===============

\begin{lstlisting}[
	style=normal,
	caption={Дерепликация PTE-таблиц},
	label={lst:dereplicate_pgtables_pte_range}
]
/* mm/numa_user_replication.c */

static void dereplicate_pgtables_pte_range(struct mmu_gather *tlb,
	pmd_t *pmd, unsigned long addr) {
	unsigned long offset;
	struct ptdesc *curr, *tmp;
	pmd_t *curr_pmd;
	pte_t *curr_pte;
	spinlock_t *lock;
	/* master-table is the only one that will remain */
	pgtable_t ptetable = ptdesc_page(
		page_ptdesc(pmd_page(*pmd))->master_table);

	/* nothing need to do if PTE table is not replicated */
	if (!numa_pgtable_replicated(page_to_virt(ptetable)))
		return;

	/* Populate PMD-entries to point to the master-table.*/

	lock = pmd_lock(tlb->mm, pmd); /* probably unnecessary */
	pmd_populate(tlb->mm, pmd, ptetable);
	for_each_pgtable_replica(curr, curr_pmd, pmd, offset)
		pmd_populate(tlb->mm, curr_pmd, ptetable);
	spin_unlock(lock);

	/* Release unneeded tables. */
	for_each_pgtable_replica(curr, tmp, curr_pte,
		page_to_virt(ptetable), offset) {
		/* destroy the cyclic linked list */
		cleanup_pte_list(ptdesc_page(curr));
		/* release the table (PG_replicated flag unsetting here) */
		pte_free_tlb(tlb, ptdesc_page(curr), addr);
		mm_dec_nr_ptes(tlb->mm);
	}

	/* destroy the cyclic linked list and
	 * unset PG_replicated flag for master-table */
	cleanup_pte_list(ptetable);
	ClearPageReplicated(ptetable);
}
\end{lstlisting}
\vspace{0.5cm}

%===============lst:data_derepl===============

\begin{lstlisting}[
	style=diff,
	caption={Дерепликация данных},
	label={lst:data_derepl}
]
 /* mm/mprotect.c */

 int
 mprotect_fixup(struct vma_iterator *vmi, struct mmu_gather *tlb,
	 struct vm_area_struct *vma, struct vm_area_struct **pprev,
	 unsigned long start, unsigned long end, unsigned long newflags)
 {
	...
<+>+#ifdef CONFIG_USER_REPLICATION</+>
<+>+	 if (!(newflags & VM_REPLICA_COMMIT) && vma_has_replicas(vma)) {</+>
<+>+		 /* We should take vma write lock before phys_deduplicate()</+>
<+>+		  * to prevent concurents to phys_duplicate() the vma's</+>
<+>+		  * range again using vma read lock.</+>
<+>+		  */</+>
<+>+		 vma_start_write(vma);</+>
<+>+		 error = phys_deduplicate(vma, start, end - start);</+>
<+>+		 if (error)</+>
<+>+			 goto fail;</+>
<+>+	 }</+>
<+>+#endif</+>
	 ...
	 vma_start_write(vma);
	 vm_flags_reset(vma, newflags);
	 ...
 }
\end{lstlisting}
